\section{Reproducibility}
\label{app:repro}

\paragraph{Software and data.} Simulation: robosuite \rf{ver.robosuite} (commit
\texttt{\rf{ver.robosuite.commit}}), MuJoCo \rf{ver.mujoco}; learning: robomimic
\rf{ver.robomimic} (commit \texttt{\rf{ver.robomimic.commit}}), PyTorch \rf{ver.torch}, NumPy
\rf{ver.numpy}, diffusers \rf{ver.diffusers} (scheduler reference only); Python
\rf{ver.python}. All versions are pinned in the repository's \file{uv.lock}. Every reported
rollout ran in the container image \code{sha256:}\rf{plat.image}. The demonstrations are the
Hugging Face \code{robomimic/robomimic_datasets} files at revision \rf{data.hf.revision}, with
the SHA-256 prefixes of Section~\ref{sec:setup}, released under the \rf{data.license} license
according to the dataset card; \code{make download} fetches and verifies them, and they are
not redistributed. The RoboForge code is released under the license Apache-2.0; robosuite,
robomimic and MuJoCo keep their own licenses.

\paragraph{Splits and initial states.} The training and validation episode IDs are in
\file{data/splits/lift_ph.json} and \file{data/splits/can_ph.json} (Table~\ref{tab:sources}).
Each evaluation suite is a manifest with per-state hashes under
\file{data/eval_states/<task>/<suite>/}; the held-out manifests' content SHA-256 prefixes are
\texttt{\rf{l1.lift.heldout.manifest}} (Lift) and \texttt{\rf{l1.can.heldout.manifest}} (Can), as
recorded in the Level~1 freeze. The robustness sets are derived deterministically from the
base starts by \file{configs/eval/robustness.yaml}.

\paragraph{Frozen results.} The Level~1 held-out analysis ran at commit \texttt{\rf{l1.git}},
the robustness analysis at \texttt{\rf{rob.git}}, the RGB analysis at \texttt{\rf{rgb.git}} and the
Level~2 analysis at \texttt{\rf{l2.git}}. Each study's freeze document, committed before its
first governed run, lists the SHA-256 of every code and configuration file it depends on; a
repository check compares those hashes with the working tree. Table~\ref{tab:sources} lists
every frozen file this report reads, with its SHA-256 prefix. The evaluation summaries that a
results file cites are read only after their hash matches the one recorded in that file.

\paragraph{Regenerating this report.} Every table, figure and number in the text is generated
from those files by \code{scripts/paper_assets.py}; the prose uses only generated macros. The
generator never reads the interactive viewer's summaries or the training-run directories.
\begin{itemize}
\item \code{make paper-assets} regenerates \file{paper/generated/}, including an index with
  the SHA-256 of every output and of every source it read.
\item \code{make paper-check} confirms that the committed generated files equal a fresh build
  byte for byte and lints the sections for typed numbers, placeholders and undefined macros.
\item \code{make paper-bibcheck} resolves every DOI and arXiv identifier of the bibliography
  and compares titles and first authors (it needs network access).
\item \code{make paper} builds \file{paper/main.pdf} with Tectonic from its cached bundle,
  without network access, and with a fixed source date so that the PDF is byte-identical
  across builds.
\end{itemize}

\paragraph{Commands and budgets.} Training runs through \code{make train} and the suite
orchestrator \code{make suite} with the committed plans
(\file{configs/experiments/level1_core_lift.yaml} and \file{level1_core_can.yaml}): every
run \rf{train.steps} steps, a checkpoint every \rf{train.ckpt} steps, every checkpoint on the
\rf{proto.nstates} validation starts. Every rollout is capped at the \rf{proto.horizon}-step
horizon. The held-out tests and their analyses (\code{make eval-final},
\code{make analyze-final}, \code{make rgb-heldout}, \code{make pc-placement},
\code{make pc-heldout}, \code{make level2-analyze}) each ran once, after the freeze that governs it was
committed, and must never be rerun. Each paid GPU launch was limited to one
hour and ran under a committed approval with a dollar cap (Table~\ref{tab:spend}).

\paragraph{What is not redistributed.} Trained checkpoints, rollout records and per-episode
arrays live in the training-run directories, which are not part of the repository. The
committed evaluation summaries, selections and analyses suffice to regenerate every number
here, but not to re-run a policy without retraining it. Retraining follows the committed
configurations; Level~1 runs train on a CPU, and the RGB and Level~2 runs need one NVIDIA L4
GPU.

\begin{table}[p]
\centering\scriptsize
\caption{Frozen sources read by the report generator, with their SHA-256 prefixes.}
\label{tab:sources}
% Generated by scripts/paper_assets.py (DECISIONS D-094). Do not edit.
% Sources (path sha256):
%   configs/models/diffusion_unet.yaml ca181c22a4af5499a67066fe2d78c06aeff5a889212ad3cceed7265fc8228bbf
%   configs/models/pointcloud/diffusion_dp3_lift.yaml b56f515dfd861568c27d5c89ebaa9878bb87452a634daebf776f58174bfe802a
%   reports/compute/cost_ledger.jsonl a4f92b9f1f4ae3708bacea7a34946b3a51705205b1fc4e4a2666096f3ba40578
%   reports/data_audit.json 1812ad718b506247ab10361d760652b6db1286f5758a05bfb847932695dbaa21
%   data/manifests/robomimic_v1.5_ph.json 34987db82f9a60a0da017ca5169eeae2214fd590505786063f59b85226f89188
%   docs/FINAL_TEST_FREEZE.md 6a3c6874c4bc1857efe46ccb05c0a1fa68247671fc9a96f5f54ce079aa9153e9
%   reports/session7/results/results.json b8b097c38765c588711615917e3cfd5be6847e63083ddebe92f0f1488672408c
%   reports/session7/robustness/analysis.json 8cca66a8a8a8cecf0b734355167936e81983c821c625a23456523365409c9451
%   reports/session7/results/taxonomy.json 7a71518171c8303cc00e357626ba66c64902b3d0e179ce82ba41504c2d6516b8
%   reports/session9/level2_results/ERRATA.md 280974a994408a6f7d3683788561f0f6c05de90f4eccf8d7b1636dd90fabca31
%   docs/LEVEL2_EXPERIMENT_FREEZE.md ffe5a03e98714465acf8655c281dea653197835f4e218cfb97a1ba6173589d81
%   reports/session9/level2_results/results.json 1d33c975e49df907a6543b17e0cf36d047f7b5ce11aa69a355313e146c22571e
%   logs/session9/pc_train/launch_records/train_gpu_pc_20261109T130430.result.json 0ea097359ba49b375e3ce53199435b771e9c3253e1e039baa62f71a003e2a737
%   reports/session4/param_counts.json 0ec7b26a0a329737cc0b18a968f5276d7130a5f66614305985505a16532486cb
%   configs/experiments/level1_core_can.yaml 881b4fb7fd72a33866ad4924db3b72cbe4270a212b03a8cc958efc6ba8a89e3c
%   configs/experiments/level2_pointcloud.yaml d91a546e1774779b27a16edbd5f33b0a2d4c11e6d68c71c13e2dff7e0a3e5c9d
%   configs/experiments/level1_core_lift.yaml d3de6cbcb14bcda15023a6526c6e3ba8cafb0b37d8692c5f8b6fcba871cfad3f
%   configs/eval/platform.yaml a152eb41157f4b4dc6f9f1767c79d2e32c6df9b9a351430ca1dfe7af24eb2fd5
%   docs/RGB_EXPERIMENT_FREEZE.md 47e7cb8ee70743dbcc6a0295a92a7462a2309870aebb3eda19060ca0e0ac385f
%   reports/session8/rgb_results/results.json 57663b81aafc70e9c1d850322b6bb1b894f7dad0d082976b1f546475a601390e
%   configs/eval/robustness.yaml cfa933cac8719e329c8159d2ccfa18ba2fcd58cc412b9ef65707e3e918a11d1f
%   data/splits/can_ph.json 754967ce665fca0084d3657e8e2d511db43d35f218d517c0bd15c93c17e75115
%   data/splits/lift_ph.json e36d89afb6d0704aa8536e23696fe30f894243ed80030e568f9a0ba3fd3a5273
%   reports/session7/suite/level1_core_can.json 213bb2c00de1a6f6631d431cb6114d2c7776ba75fc47cbc59c963ebd9adb8629
%   reports/session6/suite/level1_core_lift.json 45c28837c1fd507946856f8acc3d9a376da4ce86de1d5222ea0a08b82c757faa
%   uv.lock 327eaa2ea05a4e8fa46d8ed717ae4c7f5bcfa9a4069edac31b778d6b42ac5ca9
\begin{tabular}{p{0.42\linewidth}lp{0.36\linewidth}}
\toprule
Source & SHA-256 (prefix) & Content \\
\midrule
\texttt{reports/\allowbreak{}session7/\allowbreak{}results/\allowbreak{}results.json} & \texttt{b8b097c38765c588} & Level 1 held-out analysis (frozen, run once) \\
\texttt{reports/\allowbreak{}session7/\allowbreak{}results/\allowbreak{}taxonomy.json} & \texttt{7a71518171c8303c} & Level 1 held-out failure taxonomy per episode \\
\texttt{reports/\allowbreak{}session7/\allowbreak{}robustness/\allowbreak{}analysis.json} & \texttt{8cca66a8a8a8cecf} & Level 1 robustness and ablation analysis (validation regime) \\
\texttt{reports/\allowbreak{}session8/\allowbreak{}rgb\_results/\allowbreak{}results.json} & \texttt{57663b81aafc70e9} & RGB study held-out analysis (frozen, run once) \\
\texttt{reports/\allowbreak{}session9/\allowbreak{}level2\_results/\allowbreak{}results.json} & \texttt{1d33c975e49df907} & Level 2 analysis, placement (H5) and held-out (frozen, run once) \\
\texttt{reports/\allowbreak{}session9/\allowbreak{}level2\_results/\allowbreak{}ERRATA.md} & \texttt{280974a994408a6f} & Level 2 erratum (KNOWN\_ISSUES KI-049), cited beside the row it corrects \\
\texttt{logs/\allowbreak{}session9/\allowbreak{}pc\_train/\allowbreak{}launch\_records/\allowbreak{}train\_gpu\_pc\_20261109T130430.result.json} & \texttt{0ea097359ba49b37} & Level 2 seed 0 completing launch record (KI-049 correction) \\
\texttt{reports/\allowbreak{}compute/\allowbreak{}cost\_ledger.jsonl} & \texttt{a4f92b9f1f4ae370} & Machine-readable cost ledger (launch, result and billing events) \\
\texttt{docs/\allowbreak{}FINAL\_TEST\_FREEZE.md} & \texttt{6a3c6874c4bc1857} & Level 1 held-out freeze (machine-readable block) \\
\texttt{docs/\allowbreak{}RGB\_EXPERIMENT\_FREEZE.md} & \texttt{47e7cb8ee70743db} & RGB study freeze \\
\texttt{docs/\allowbreak{}LEVEL2\_EXPERIMENT\_FREEZE.md} & \texttt{ffe5a03e98714465} & Level 2 freeze \\
\texttt{data/\allowbreak{}manifests/\allowbreak{}robomimic\_v1.5\_ph.json} & \texttt{34987db82f9a60a0} & Dataset manifest (robomimic v1.5 PH, Hugging Face revision, SHA-256) \\
\texttt{data/\allowbreak{}splits/\allowbreak{}lift\_ph.json} & \texttt{e36d89afb6d0704a} & Lift demonstration split \\
\texttt{data/\allowbreak{}splits/\allowbreak{}can\_ph.json} & \texttt{754967ce665fca00} & Can demonstration split \\
\texttt{reports/\allowbreak{}data\_audit.json} & \texttt{1812ad718b506247} & Data audit (episode lengths, control period, action statistics) \\
\texttt{reports/\allowbreak{}session4/\allowbreak{}param\_counts.json} & \texttt{0ec7b26a0a329737} & Exact parameter counts of the Level 1 families \\
\texttt{configs/\allowbreak{}eval/\allowbreak{}robustness.yaml} & \texttt{cfa933cac8719e32} & Preregistered robustness sets (RESEARCH\_PROTOCOL §14) \\
\texttt{configs/\allowbreak{}eval/\allowbreak{}platform.yaml} & \texttt{a152eb41157f4b4d} & Canonical evaluation platform \\
\texttt{configs/\allowbreak{}experiments/\allowbreak{}level1\_core\_lift.yaml} & \texttt{d3de6cbcb14bcda1} & Level 1 Lift suite plan (budget, tuning trials) \\
\texttt{configs/\allowbreak{}experiments/\allowbreak{}level1\_core\_can.yaml} & \texttt{881b4fb7fd72a338} & Level 1 Can suite plan (budget, tuning trials) \\
\texttt{configs/\allowbreak{}experiments/\allowbreak{}level2\_pointcloud.yaml} & \texttt{d91a546e1774779b} & Level 2 plan \\
\texttt{configs/\allowbreak{}models/\allowbreak{}diffusion\_unet.yaml} & \texttt{ca181c22a4af5499} & Level 1 diffusion configuration (base; trials override the optimizer only) \\
\texttt{configs/\allowbreak{}models/\allowbreak{}pointcloud/\allowbreak{}diffusion\_dp3\_lift.yaml} & \texttt{b56f515dfd861568} & Level 2 point-cloud arm configuration \\
\texttt{reports/\allowbreak{}session6/\allowbreak{}suite/\allowbreak{}level1\_core\_lift.json} & \texttt{45c28837c1fd5079} & Level 1 Lift suite summary (validation selections, development numbers) \\
\texttt{reports/\allowbreak{}session7/\allowbreak{}suite/\allowbreak{}level1\_core\_can.json} & \texttt{213bb2c00de1a6f6} & Level 1 Can suite summary (validation selections, development numbers) \\
\texttt{uv.lock} & \texttt{327eaa2ea05a4e8f} & Locked Python environment (package versions) \\
\texttt{reports/\allowbreak{}session10/\allowbreak{}observatory/\allowbreak{}verification/\allowbreak{}lift-\allowbreak{}diffusion\_unet-\allowbreak{}s0-\allowbreak{}lift\_final\_test\_0000\_t021\_frontview\_rgb.png} & \texttt{b6a165f72de1fb92} & Lift, diffusion seed 0, held-out start 0, frame 21: MuJoCo frontview render of the recorded state (Session 10 viewer-verification reference) \\
\texttt{reports/\allowbreak{}session10/\allowbreak{}observatory/\allowbreak{}verification/\allowbreak{}can-\allowbreak{}diffusion\_unet-\allowbreak{}s0-\allowbreak{}can\_final\_test\_0000\_t060\_frontview\_rgb.png} & \texttt{149b860b17a2eca4} & Can, diffusion seed 0, held-out start 0, frame 60: MuJoCo frontview render of the recorded state (Session 10 viewer-verification reference) \\
\texttt{reports/\allowbreak{}session9/\allowbreak{}pc\_feasibility/\allowbreak{}linux-\allowbreak{}aarch64-\allowbreak{}container-\allowbreak{}1f685638f28c/\allowbreak{}lift\_overlay\_demo\_0\_t029.png} & \texttt{fe163485aaccd6a2} & Lift demo 0, t = 29: the 512 sampled cloud points projected on the frontview and agentview frames (Session 9 feasibility, canonical container) \\
\texttt{reports/\allowbreak{}session7/\allowbreak{}heldout/\allowbreak{}} (26 files) & per file & Level 1 held-out evaluation summaries, each read only after its SHA-256 equals the one recorded in the Level 1 results file (D-096) \\
\bottomrule
\end{tabular}

\end{table}
